Denna kandidatuppsats inriktar sig mot grundläggande koncept inom algebraisk geometri, eliminering av polynoma system samt Gröbner baser. Algebraisk geometri snuddades vid genom att introducera affina rum och varieteter och hoppar sedan vidare till algebra genom att utforska Gröbner baser. Utforskningen av dessa baser utförs med konstruktionen av en algoritm som kan beräkna en Gröbner bas som slutmål. För att uppnå detta diskuteras monomialideal och polynomideal samt användbara egenskaper. Därefter bevisas den grundläggande kärnan till Gröbner bas algoritmen, Buchberger kriterium. Detta kriterium används för att programmera Buchbergers algoritm, alltså skapandet av en Gröbner bas, i datorsalgebrapråket Macaulay2. Envariabel och flervariabel polynomdivision samt lösningar till polynoma system implementerades också i Macaulay2. Resultaten verifierades med enhetstester. Uppsatsen lägger en stark grund till framtida utforskning av algebraisk geometri och algoritmkonstruktion.   